\documentclass[9pt]{beamer-control}
\usepackage{beamer-control-workshop}
\begin{document}
\CONCEPT[6]{Week 6: Frequency Domain Analysis}

\begin{frame}
\frametitle{Introduction}
In this workshop, we will explore frequency domain analysis techniques including Bode plots and Nyquist plots.
\end{frame}

\SUBCONCEPT{More on the frequency domain}

\begin{frame}[fragile]
\frametitle{Nyquist plots}
Last workshop we saw how to graph the Bode plots for a given transfer function, now we can include Nyquist plots. The command for Nyquist plots is not hard to guess,
\includematlab{workshop6.m}{part1} 

Use Matlab to define the following transfer functions:
\[ G_1(s) = \frac{1}{s+10},  \quad  G_2(s) = \frac{1000}{(s+10)^3} \]
\[ G_3(s) = \frac{s/100+1}{(s/10+1)(s/1000+1)} , \quad  G_4(s) = \frac{s}{s+10} \]
\[ G_5(s) = \frac{1}{s(s+10)} \]
Plot the Bode and Nyquist plots for $G_{1:5}$ and see if you can spot typical features. Change poles, zeros, gains, and otherwise explore the landscape of $s$-domain functions until you feel like you have some intuition behind the relationships connecting the mathematical function and the graphical representation.
\end{frame}


\begin{frame}[fragile]
	\frametitle{A particular example}
	Let us look closer at a particular example, the transfer function
	\[L(s) =\frac{k}{(s+1)^3}\]
	Here our gain is given by the parameter $k$, let us first consider $k=1$.\\[10pt]
	You may define the closed loop transfer function using the $\texttt{feedback()}$ command
	\includematlab{workshop6.m}{part2} 
	
	Investigate both the step response (of the closed loop system) and Nyquist plot (of the open loop system) for values of $k\in[1,10]$. \\[10pt]
	You should notice that for a sufficiently large gain  value, the closed loop system becomes unstable, and this corresponds to encirclement of the critical point $s=-1$ on the Nyquist plot.
\end{frame}


\SUBCONCEPT{Stability margins}
\begin{frame}[fragile]
	\frametitle{Calculating margins}
	A notion of how close we are from instability can be gained from stability margins. The gain and phase margin may calculated using a Nyquist plot and tell us how much gain or phase can be added to the system before becoming unstable.\\[10pt]
	
	Using the same system as before and setting $k=5$, graph the Nyquist plot of the open loop system. If you right-click on the graph and click Characteristics > All Stability Margins, you can see these margins directly on the Nyquist plot. The same applies to a Bode plot (try this out!). The values of these margins can also be found using 
	\includematlab{workshop6.m}{part3} 
	Play around with different transfer functions, such as $G_{1:5}$ shown earlier in this Workshop.
	
\end{frame}




\end{document}
